% !TEX root = tarasov-qasp-zh.tex


\chapter*{答案}
\phantomsection
\addcontentsline{toc}{chapter}{答案}
%\markboth{答案}
\label{answers}

\begin{fullwidth}
\begin{multicols}{2}
\begin{enumerate}[label=\arabic*.]
\item \SI{20}{\metre}；\SI{1}{\second}；$v_{A} = \SI[per-mode=symbol]{10.2}{\metre\per\second}$；$v_{B} = \SI[per-mode=symbol]{10.6}{\metre\per\second}.$
\item $v_{0}= \SI[per-mode=symbol]{11.3}{\metre\per\second}; \,\, x = \SI{4}{\metre}; \,\, y =  \SI{0.8}{\metre}; \,\, t = \SI{0.5}{\second};$\\
$ v_{A}=\SI[per-mode=symbol]{9.4}{\metre\per\second}; \,\,v_{B}=\SI[per-mode=symbol]{15.2}{\metre\per\second}$。
\item
\begin{enumerate}[label=(\arabic*)]
\item $t \sqrt{v_{1}^{2} + v_{2}^{2} + 2v_{1}v_{2} \cos (\alpha_{1} + \alpha_{2})}$；
\item $t \sqrt{v_{1}^{2} + v_{2}^{2} - 2v_{1}v_{2} \sin \alpha_{1} \sin \alpha_{2} }$。
\end{enumerate}
\item $\dfrac{3}{2} \sqrt{2 \left( \dfrac{H-h}{g} \right)} + \sqrt{\dfrac{H+3h}{2g}}$。
\item $\cot \alpha = \left( \dfrac{P-4F}{4P} \right)$。
\item $ 2 \left( \dfrac{v_{0}^{2}}{g} \right) \dfrac{\sin \alpha}{\cos^{2} \alpha}$。
\item \SI[per-mode=symbol]{13.8}{\metre\per\second}。
\item \SI[per-mode=symbol]{37.2}{\metre\per\second}；\SI{1280}{\joule}。
\item \SI[per-mode=symbol]{2.6}{\metre\per\second\squared}；\SI{42}{\newton}；\SI{8.5}{\newton}。
\item \SI[per-mode=symbol]{3.3}{\metre\per\second\squared}；\SI{13}{\newton}。
\item \SI[per-mode=symbol]{3.5}{\metre\per\second\squared}；\SI{33.6}{\newton}；\SI{50.4}{\newton}。
\item \SI[per-mode=symbol]{6.9}{\metre\per\second\squared}；\SI{8.8}{\newton}；\SI{16.2}{\newton}；\SI{1.5}{\newton}。
\addtocounter{enumi}{2}
% TODO-FOR-CHECK（结构核对，非疑误）：底本答案表缺第 13、14 题（第 7 章两道
% fig-32 讨论题）的答案，编号在此跳至 15。已核对全书习题连续编号共 82 题：第 8 章
% 图 39/40 图注自证第 16、19 题，第 25 章题面自证第 45 题，末号 82 与第 33 章末题
% 吻合。故 \addtocounter{enumi}{2} 系对印刷编号的忠实复现，非录入错位，照抄。
\item 0.45。
\item 7:4:1。
\item $\sqrt{5gR}$。
\item $\dfrac{81 \pi}{8 GT^{2}}$。
\item  $h=R \left(1- \dfrac{g}{\omega^{2} R} \right);\,\, F=m\omega^{2} R$。
\item 1.5R。
\item \SI[per-mode=symbol]{120}{\kilo\gram\per\metre\cubed}。
\item \SI{3900}{\joule}。
\item 0.27。
\item 0.5。
\item
$F =mg \left( \dfrac{H}{h} + 1 \right) - V g \rho_{w}; $\\
$h_{1} = 2h \left( \dfrac{V \rho_{w}}{m} - 1 \right) - H.$

\item \SI[per-mode=symbol]{7.5}{\kilo\metre\per\hour}；\SI{4.65}{\metre}。
\item $h = \dfrac{v_{0}^{2}}{2g} \left( \dfrac{m}{M}\right)^{2} \dfrac{\cos^{2} \alpha \sin \alpha}{\sin \alpha + k \cos \alpha}$。
\item
\begin{enumerate}[label=(\arabic*)]
\item $h_{1} = \dfrac{2}{9g} \left( 4gl \sin^{2} \dfrac{\alpha}{2} + v_{0}^{2}\right);$\\
$h_{2} = \dfrac{1}{18g} \left( 4gl \sin^{2} \dfrac{\alpha}{2} + v_{0}^{2}\right);$
\item $h = \dfrac{1}{18g} \left( 4gl \sin^{2} \dfrac{\alpha}{2} + v_{0}^{2}\right)$。
\end{enumerate}
\item $v_{min} = \left( \dfrac{m+M}{m} \right) \sqrt{5gl}$。
\item $H \left( \dfrac{M-m}{M+m} \right)$。
\item $\dfrac{HM}{4M+3m}$。
\item $\dfrac{3}{4} \sqrt{\dfrac{3}{2}}$。
\item 4。
\item \SI[per-mode=symbol]{0.43}{\metre\per\second}。
\item \SI[per-mode=symbol]{27.4}{\metre\per\second\squared}：加速度方向竖直向上。
\item \SI{1.28}{\newton}；\SI{1.28}{\newton}；\SI{0.62}{\newton}；\SI{1.56}{\newton}。
\item 在圆盘中心右侧 $\dfrac{3}{22}R$ 处。
\item $0.05 \dfrac{V}{S}$。
\item \SI{11.3}{\centi\metre}；\SI{13.4}{\milli\gram}。
% TODO-FOR-CHECK: 底本此处原文为 ``\SI{11.3}{\centi\metre} em''，"em" 疑为衍文
% （无此单位；按题设等温过程验算，管内空气柱 11.3 cm 对应空气质量约 13.3 mg，
% 与后值 13.4 mg 吻合），中文版径删，待考原书。
\item 下降 \SI{3}{\centi\metre}；\SI{15.4}{\milli\gram}。
\item \SI{59}{\gram}。
\item
\begin{enumerate}[label=(\arabic*)]
\item \SI{138}{\joule}；
\item  \SI{171}{\joule}。
\end{enumerate}
\item 伸长 \SI{1.5}{\centi\metre}；\SI{21.5}{\milli\gram}。
\item \SI{735}{\gram}；不会结露；0.58\%。
\item \SI{3d-8}{\second}；\SI[per-mode=symbol]{5d8}{\metre\per\second}。
% TODO-FOR-CHECK: 原值 5×10^8 m/s 疑误。第 45 题（第 25 章）板长 15 cm、
% E = 500 V/m、板间距 4 cm，算得 t = 3×10^-8 s（与前值吻合）、
% v_min = 0.15 m / (3×10^-8 s) = 5×10^6 m/s；且 5×10^8 m/s 超光速，对静电场中
% 加速的电子不可能。疑为 5d6 之误。照译原值，待考原书。
\item \SI[per-mode=symbol]{147}{\volt\per\metre}。
\item  $\dfrac{(mg+ Eq)}{\cos \alpha}; \,\, \dfrac{(mg+Eq) l \sin \alpha \tan \alpha}{2}.$
\item $\dfrac{(m_{1} - m_{2})g+E(q_{1}-q_{2})}{(m_{1} + m_{2})};$\\
$\dfrac{m_{1} m_{2}g + (m_{2}q_{1} +  m_{1}q_{2} )E}{(m_{1} + m_{2})}.$
\item $\sqrt{5\left( mg +Eq \right) \left( \dfrac{l}{m} \right)}$。
\item 1.83 q。
\item $\sqrt{\dfrac{g}{l \cos \alpha} - \dfrac{q^{2}}{ml^{3} sin^{3} \alpha}};\,\, \left( \dfrac{mg}{\cos \alpha} \right)$。
\item $\sqrt{5gl - \dfrac{q^{2}}{ml}}$。
\item \SI{0.2}{\ampere}。
\item \SI{1}{\ampere}。
\item \SI{0.16}{\ohm}。
\item \SI{9.9}{\volt}；1\%。
\item \SI{0.196}{\ampere}；1.96\%。
\item \SI{6d-6}{\coulomb}。
\item $\dfrac{\Ea C R_{1}R_{2}}{R_{1}r+2R_{2}r+2R_{1}R_{2}}$。
\item $\dfrac{rdmv_{0}^{2} \tan \alpha}{\Ea ql - dmv_{0}^{2} \tan \alpha}$。
\item $\dfrac{rdmv_{0}^{2}}{\Ea ql - dmv_{0}^{2}}$。
\item $\arctan \left( \dfrac{3 \Ea ql }{mdv_{0}^{2}} \right)$。
\item \SI{3.75}{\volt}；\SI{0.25}{\ampere}。
\item $\dfrac{2}{3} R$。
\item $\dfrac{4}{5} R$。
\item $\dfrac{3}{4} R$。
\item $\dfrac{11}{20} R$。
\item \SI{60}{\ohm}；\SI{70}{\watt}。
\item 减小 28.6\%。
\item $\dfrac{R}{8}$；89\%；83\%。
\item \SI{800}{\gram}。
\item 有 \SI{100}{\gram} 水化为蒸汽；\SI{21}{\minute}。
\item \SI{10.8}{\centi\metre}。
\item \SI{7.4}{\metre}。
\item \ang{56}；\SI{2.3}{\centi\metre}。
\item $\dfrac{d}{n} \left( 1 - \dfrac{1}{\sqrt{n^{2}+1}}\right).$
\item 是发散透镜；$d= \SI{15}{\centi\metre}$；像将远离透镜移动 \SI{6}{\centi\metre}；$k = 0.4$。
\item $f = \SI{37.5}{\centi\metre}$；像将变为实像；$f_{1} = \SI{150}{\centi\metre}$。
\item 在会聚透镜右侧 \SI{100}{\centi\metre} 处。
\item 位于中间那块透镜的中心。
\item $f=\SI{100}{\centi\metre}; \,\, k=4$。
\item $f = \SI{6.3}{\centi\metre}; \,\,k = 0.63$。
\end{enumerate}
\end{multicols}
\rule{14cm}{1pt}
\vspace{2cm}
\begin{center}
{\Huge\bfseries 全书完}
\end{center}


\end{fullwidth}

%\newpage
%\thispagestyle{empty}
%\begin{fullwidth}
%\vfill
%\vspace*{10cm}
%\begin{center}
%\allcaps{FINIS}
%\textls[500]{FINIS}
%\end{center}
%\end{fullwidth}
\cleardoublepage

% NEWTERMS: converging lens | 会聚透镜 | 承 ch-31/32 定名（人教版初中称凸透镜）；本章答案 79
% NEWTERMS: diverging lens | 发散透镜 | 承 ch-31/32 定名（人教版初中称凹透镜）；本章答案 77
% NEWTERMS: real image | 实像 | 承 ch-32 定名；本章答案 78“像将变为实像”
% NEWTERMS: steam | 蒸汽 | 承 ch-30 习题 5 题面定名“化为蒸汽”；本章答案 72
% NEWTERMS: dew | 露 | 承 ch-21 习题 6 题面定名“结露”；本章答案 44“不会结露”
